% Source: https://artofproblemsolving.com/community/c3h3064003_my_math_story_for_325_year_anniversary
% Stable identity: aops-topic-3064003-q15
Put $A=(0,0)$, $B=(1,0)$, $C=(1,1)$ and $D=(0,1)$. The four lines are $y=\dfrac{x}{2}$, $y=2-2x$, $y=1-\dfrac{x}{2}$ and $y=2x$. Their successive vertices around the bounded quadrilateral are $A$, $P=\left(\dfrac{4}{5},\dfrac{2}{5}\right)$, $Q=\left(\dfrac{2}{3},\dfrac{2}{3}\right)$ and $R=\left(\dfrac{2}{5},\dfrac{4}{5}\right)$. The shoelace formula gives area $\dfrac{1}{2}\left[\left(\dfrac{8}{15}-\dfrac{4}{15}\right)+\left(\dfrac{8}{15}-\dfrac{4}{15}\right)\right]=\boxed{\dfrac{4}{15}}$.
\begin{asy}
usepackage("amsmath");
size(7cm);
pair A=(0,0),B=(1,0),C=(1,1),D=(0,1),M=(1,.5),Npt=(.5,1),P=(.8,.4),Q=(2.0/3,2.0/3),R=(.4,.8);fill(A--P--Q--R--cycle,lightgray);draw(A--B--C--D--cycle);draw(A--M);draw(B--Npt);draw(M--D);draw(A--Npt);label("$A$",A,SW);label("$B$",B,SE);label("$C$",C,NE);label("$D$",D,NW);label("$M$",M,E);label("$N$",Npt,N);label("$P$",P,SE);label("$Q$",Q,NE);label("$R$",R,NW);
\end{asy}
